\documentclass[10pt]{article}

% ---------------------------------------------------------------------------
% First full draft. Written BLIND to the Phase 1 outcomes: every confirmatory number is a
% \RESULT{key} placeholder, filled in after analysis/analyze.py has run on data/phase1.
% Build: tectonic main.tex   (bibliography: refs.bib, run automatically)
% Placeholder keys are also written to the log as "RESULTKEY: <key>".
% ---------------------------------------------------------------------------
\usepackage[utf8]{inputenc}
\usepackage[T1]{fontenc}
\usepackage{amsmath,amssymb}
\usepackage{newtxtext,newtxmath}
\usepackage[a4paper,margin=1in]{geometry}
\setlength{\emergencystretch}{3em}
\usepackage[htt]{hyphenat}
\usepackage{microtype}
\usepackage{booktabs}
\usepackage{tabularx}
\usepackage{multirow}
\usepackage{graphicx}
\usepackage{xcolor}
\usepackage{enumitem}
\usepackage{tikz}
\usetikzlibrary{positioning,arrows.meta,fit}
\definecolor{linknavy}{rgb}{0.13,0.29,0.53}
\definecolor{resred}{rgb}{0.70,0.05,0.05}
\usepackage{titlesec}
\usepackage{fancyhdr}
\usepackage[font=small,labelfont=bf,skip=6pt]{caption}
\usepackage[numbers,sort&compress]{natbib}
\setlength{\bibsep}{1.5pt plus 0.5pt}
\renewcommand{\bibfont}{\small}
\usepackage[colorlinks=true,allcolors=linknavy]{hyperref}

\titleformat{\section}{\normalfont\large\bfseries}{\thesection}{0.6em}{}
\titleformat{\subsection}{\normalfont\normalsize\bfseries}{\thesubsection}{0.6em}{}
\titlespacing*{\section}{0pt}{1.5ex plus 1ex minus .2ex}{0.8ex}
\titlespacing*{\subsection}{0pt}{1.2ex plus 1ex minus .2ex}{0.5ex}

\pagestyle{fancy}\fancyhf{}
\renewcommand{\headrulewidth}{0.3pt}
\fancyhead[L]{\small\itshape Architecture and AI Coding Agents: a Preregistered Experiment}
\fancyhead[R]{\small\itshape Draft --- results blinded}
\fancyfoot[C]{\small\thepage}
\fancypagestyle{plain}{\fancyhf{}\fancyfoot[C]{\small\thepage}\renewcommand{\headrulewidth}{0pt}}

\renewenvironment{abstract}{%
  \begin{center}\rule{\textwidth}{0.4pt}\end{center}\vspace{-0.9em}
  \begin{quotation}\small\noindent\textbf{Abstract.}\enspace}
  {\end{quotation}\vspace{-0.6em}\begin{center}\rule{\textwidth}{0.4pt}\end{center}\vspace{0.4em}}

% ---------------------------------------------------------------- placeholders
% \RESULT{key}: a confirmatory or descriptive number from Phase 1, not yet known.
\newcommand{\RESULT}[1]{\typeout{RESULTKEY: #1}\textcolor{resred}{\textbf{[#1]}}}
% \TODO{text}: an author note that must be resolved before submission.
\newcommand{\TODO}[1]{\textcolor{orange!70!black}{\textsf{\small[TODO: #1]}}}
% One row of a contrast table: estimate, 95% CI, p, Holm-adjusted p.
% #1 hypothesis, #2 contrast key, #3 row label, #4 estimate suffix
\newcommand{\crow}[4]{#3 & \RESULT{#1-#2-#4} & \RESULT{#1-#2-ci} & \RESULT{#1-#2-p} & \RESULT{#1-#2-pholm}\\}
% Figure that compiles before the figure file exists.
\newcommand{\maybefig}[2]{%
  \IfFileExists{#1}{\includegraphics[width=#2]{#1}}{%
    \fbox{\parbox[c][3.2cm][c]{#2}{\centering\small\textsf{Figure pending: \texttt{\detokenize{#1}}}\\[2pt]
    \textsf{(produced from data/phase1 after unblinding)}}}}}

\newcommand{\NONE}{\textsc{None}}
\newcommand{\LAY}{\textsc{Layered}}
\newcommand{\HEX}{\textsc{Hexagonal}}
\newcommand{\code}[1]{\texttt{#1}}

\title{\vspace{-1.3em}\bfseries\Large Does Architecture Help or Hinder AI Coding Agents?\\[0.4em]
\large\mdseries A Preregistered Experiment with Layered and Hexagonal Variants\\ of an Evolving Java Service}

\author{%
  Tomasz Pisarewski\\
  \normalsize Fide-Soft\\
  \normalsize \texttt{tomasz.pisarewski@fide-soft.com}}

\date{Draft of \today}

\begin{document}
\maketitle

\begin{abstract}
\noindent
Software architecture exists to help people understand and change code. Coding agents now make a
growing share of those changes. Two questions follow. Does a structure designed for human
cognition help or hinder an agent? And does the agent keep that structure intact as it works? We
report a preregistered controlled experiment (OSF \code{r5t3d}). We built an HL7 FHIR R4 REST
service in Java~21 and Spring Boot~3.5 in three behaviorally equivalent variants: no explicit
architecture (\NONE), layered (\LAY), and hexagonal, i.e.\ ports and adapters (\HEX). All three
pass the same 79-test black-box suite. Two open-weights models, Laguna~S~2.1 and
Qwen3-Coder-30B-A3B, run in the mini-swe-agent scaffold. Each implements the same 24 product
tickets in a fixed order, every ticket on top of its own previous changes, in five independent
runs per cell (30 trajectories, 720 attempts). Hidden acceptance tests measure success. Tokens
measure cost. After every task we measure erosion, using an architecture-agnostic index and
ArchUnit conformance. Success (H3): \RESULT{H3-abstract} (hexagonal minus none:
\RESULT{H3-hex-none-pp} percentage points, Holm-adjusted $p=$ \RESULT{H3-hex-none-pholm}).
Cost (H1): \RESULT{H1-abstract}. Erosion (H2): \RESULT{H2-abstract}. Conformance (H6):
\RESULT{H6-abstract}. Capability interaction (H4): \RESULT{H4-abstract}. The
whole experiment ran on a laptop with Docker and cost \RESULT{cost-total} in API fees (budget
USD~30). We release all variants, tasks, hidden tests, harness, transcripts, and per-task
snapshots.
\end{abstract}

\noindent\textbf{Keywords:} coding agents; software architecture; architecture erosion;
hexagonal architecture; controlled experiment; preregistration; open science.

% =====================================================================================
\section{Introduction}\label{sec:intro}

Software architecture is a tool for human cognition. Modules hide design decisions so that a
person can change one part of a system without understanding the rest~\citep{parnas1972criteria}.
Layers and ports separate concerns so that the developer knows where a change belongs and what it
may touch~\citep{buschmann1996posa,fowler2002peaa,cockburn2005hexagonal}. These structures cost
something: more files, more indirection, more mapping code. We pay that cost because it buys
comprehension and controlled change for human maintainers.

LLM-based coding agents now take on a growing share of maintenance and feature work. They read
code through a shell, not through an IDE and a mental model built over months. Two things are
unclear. First, does the indirection that helps humans help agents, or does it tax them? Our
earlier exploratory study on small Python repositories found a tax: equivalent code with ports
and adapters lowered single-task success, more so for weaker models, mainly because the logic to
change sat in opaquely named adapter files away from the entry point~\citep{pisarewski2026indirection}.
That study covered neither realistic project size, nor Java, nor sustained evolution. Second, do
agents preserve an architecture while they evolve the code? Architecture erosion, the growing gap
between intended and implemented architecture, is a classic problem of long-lived
systems~\citep{perry1992foundations,lehman1980programs,desilva2012erosion,li2022erosionmapping}.
Agents may speed it up. They optimize for local correctness, the ticket at hand and its tests,
and may drift from global architectural intent~\citep{agenticentropy2026}. A recent benchmark has agents extend their own code over checkpoints and measures erosion at
each step~\citep{slopcodebench2026}. None of these studies vary the architecture while holding
behavior constant.

We run that experiment. One FHIR R4 service exists in three behaviorally equivalent variants
(\NONE, \LAY, \HEX) that pass the same black-box suite. An agent implements 24 product tickets in
order, each on top of its own previous result, and we measure success, cost and erosion after
every task. The design, hypotheses, analysis code and sample size were frozen and registered
before any confirmatory run. We ask three research questions:

\begin{itemize}[leftmargin=3.2em,itemsep=1pt]
  \item[\textbf{RQ1}] \emph{Cost.} Does architecture change how much work an agent needs per
    task? (H1; by task type, H5)
  \item[\textbf{RQ2}] \emph{Success.} Does architecture change whether the agent completes the
    task, and does that depend on model capability? (H3; H4)
  \item[\textbf{RQ3}] \emph{Erosion.} Does an explicit architecture slow structural decay under
    agent-driven evolution, and do agents keep to the architecture's own rules? (H2; H6)
\end{itemize}

\paragraph{Contributions.}
\begin{enumerate}[itemsep=1pt]
  \item \textbf{Controlled, equivalent variants.} Three architectures of one non-trivial
    service (2.2--4.1\,kLOC), with an identical external contract, schema and dependency set,
    admitted only after a mechanical equivalence gate. Building them in parallel exposed 13
    specification ambiguities, which became binding clarifications.
  \item \textbf{Cumulative evolution with erosion tracking.} Agents work on their own earlier
    output across 24 tickets. Every task is a git snapshot, graded against all hidden tests so far
    and measured with a pinned erosion pipeline (ArchUnit, CK, PMD/CPD, SpotBugs).
  \item \textbf{Preregistration.} Hypotheses, tests, sample size and the analysis script were
    frozen (git tag \code{prereg-v1}) and registered on OSF before data collection. This paper was
    drafted before the results were known. Deviations are logged.
  \item \textbf{Open, reproducible, cheap.} Open-weights models, an offline Docker sandbox and an
    unmodified public scaffold (mini-swe-agent~2.4.6). The full experiment cost about
    USD~0.04 per attempt. All materials, transcripts and snapshots are released.
  \item \textbf{Lessons for low-budget agent experiments} (\S\ref{sec:lessons}): failure modes of
    free and paid endpoints, scaffolds and graders that we hit and fixed before registration.
\end{enumerate}

% =====================================================================================
\section{Related Work}\label{sec:related}

\paragraph{Coding-agent benchmarks.} SWE-bench~\citep{jimenez2024swebench} established
repository-level issue resolution with hidden FAIL\_TO\_PASS and PASS\_TO\_PASS tests. We follow
that grading scheme. SWE-agent~\citep{yang2024sweagent} showed that the agent--computer interface
matters. mini-swe-agent~\citep{minisweagent} reduces the scaffold to a bash-only loop and is the
reference scaffold of the SWE-bench Bash Only leaderboard~\citep{swebenchbashonly}. We use it so
that scaffold effects stay small and identical across models. Java issue-resolution benchmarks
exist~\citep{zan2024swebenchjava,zan2025multiswebench}. Like most benchmarks they evaluate
independent tasks on fixed snapshots of real repositories, so they cannot separate the effect of
architecture from everything else that differs between repositories.

\paragraph{Long-horizon evolution and degradation.} SWE-EVO~\citep{sweevo2025} poses
long-horizon evolution tasks derived from release notes. SlopCodeBench~\citep{slopcodebench2026}
is the closest design to ours: agents extend their own code over checkpoints, and erosion and
verbosity are measured at every step. It does not manipulate architecture while holding behavior
constant, and that manipulation is what this study adds. The notions of \emph{agentic
entropy}~\citep{agenticentropy2026}, where agents optimize local correctness and drift from global
architectural intent, and \emph{coherence debt} in repository-scale tasks~\citep{coherencedebt2026}
give conceptual framings for H2 and H6.

\paragraph{Code properties that affect agents.} A growing line of work holds the agent fixed and
varies the code. Trivedi and Schmitt~\citep{trivedi2026cleanliness} build minimal pairs that
differ in code cleanliness and measure success and token cost with hidden tests. They leave open
whether such effects compound over long-horizon evolution, which our cumulative design addresses
for architecture. Borg et al.~\citep{borg2026codehealth} relate file-level CodeHealth to
LLM refactoring outcomes. Ustynov~\citep{ustynov2026beyond} revisits conventions written for human
readers in light of agentic development and leaves architecture to future work. \emph{Theory of Code Space}~\citep{sapunov2026tocs}
asks whether agents can build an architectural map of a code base. Our predecessor
study~\citep{pisarewski2026indirection} asked whether indirection changes task success and traced
the effect mainly to file naming. All of these work with single tasks or single snapshots.

\paragraph{Architecture and agents.} Three recent works are closest in topic. \emph{Architecture as
Capability Equalizer}~\citep{archequalizer2026} varies how an architecture is \emph{specified} to
six models, including ArchUnit-style rules, and finds that structured specifications help weaker
models most. It does not vary the architecture of the code. A pilot study~\citep{hexdebt2025} has
three models build a hexagonal microservice and reports that, without automated architectural
linting, structural debt accumulates faster with small models. It studies one style with no
controlled alternative. In grey literature, Stoyanov~\citep{stoyanov2026hexagonal} compares a flat
and a hexagonal version of one Java system on 9 + 6 tasks and reports that the hexagonal version
took 37.8\% longer to reach acceptance. That comparison involves one system and one author and was
not preregistered. Together these works motivate H1, H2 and H6 but do not test them under control.

\paragraph{Architecture erosion.} Perry and Wolf~\citep{perry1992foundations} distinguish
\emph{erosion}, which violates the architecture, from \emph{drift}, which is insensitivity to it.
De Silva and Balasubramaniam~\citep{desilva2012erosion} survey methods to control erosion,
including conformance checking against rules, which we implement with
ArchUnit~\citep{archunit}. Li et al.\ map the erosion literature~\citep{li2022erosionmapping} and
report practitioners' views of its causes and symptoms~\citep{li2021erosionpractitioners}. Common
symptoms (dependency cycles, rising coupling, duplication, rule violations) are what our erosion
index aggregates. We measure them with established metrics: CK~\citep{chidamber1994metrics,aniche2015ck},
Martin's distance from the main sequence~\citep{martin2002agile}, PMD/CPD~\citep{pmd} and
SpotBugs~\citep{spotbugs}. To our knowledge erosion research has so far studied human developers.
Agents as the source of erosion are new.

% =====================================================================================
\section{Study Design}\label{sec:design}

This section summarizes the frozen protocol (PREREGISTRATION.md at tag \code{prereg-v1}). Phase~1,
reported here, covers the FHIR domain. A second domain (order management) is built with the same
method and will be registered separately as Phase~2, where domain generality (H7) is tested.
Figure~\ref{fig:pipeline} shows the data flow for one trajectory.

\begin{figure}[t]\centering
\begin{tikzpicture}[
  box/.style={draw, rounded corners=2pt, align=center, font=\footnotesize, minimum height=1.05cm, inner sep=3pt},
  arr/.style={-{Stealth[length=2mm]}, thick},
  node distance=0.45cm and 0.5cm]
\node[box, text width=2.25cm] (base) {Baseline variant\\ \NONE{} / \LAY{} / \HEX};
\node[box, text width=3.1cm, right=of base] (agent) {Ticket $k$ $\rightarrow$ mini-swe-agent\\ offline sandbox, $\le$60 calls};
\node[box, text width=2.0cm, right=of agent] (commit) {git commit\\ snapshot $k$};
\node[box, text width=3.6cm, right=of commit] (grade) {Grading container\\ build, start, baseline suite\\ + hidden tests of tasks $1..k$};
\node[box, text width=3.6cm, below=0.55cm of grade] (res) {\code{results.jsonl}: success,\\ regressions, tokens, calls, time};
\node[box, text width=4.4cm, below=0.55cm of agent, xshift=0.4cm] (ero) {Erosion container (post hoc):\\ ArchUnit, CK, PMD/CPD, SpotBugs\\ $\rightarrow$ \code{erosion.jsonl}};
\draw[arr] (base) -- (agent);
\draw[arr] (agent) -- (commit);
\draw[arr] (commit) -- (grade);
\draw[arr] (grade) -- (res);
\draw[arr] (commit.south) |- (ero.east);
\draw[arr, dashed] (commit.north) -- ++(0,0.45) -| node[pos=0.25, above, font=\footnotesize] {next ticket $k{+}1$ on top of snapshot $k$} (agent.north);
\end{tikzpicture}
\caption{One trajectory. The agent never sees the hidden tests. Grading and erosion measurement run
on the committed snapshot in separate offline containers.}
\label{fig:pipeline}
\end{figure}

\subsection{Application and variants}\label{sec:variants}

The application is a FHIR R4 REST facade~\citep{hl7fhirr4} over its own persistence. It covers
five resource types (\code{Patient}, \code{Practitioner}, \code{Encounter}, \code{Observation},
\code{Appointment}) with read, create, update and a defined subset of search parameters, resource
versioning (\code{meta.versionId}, \code{ETag}) and \code{OperationOutcome} errors. Every variant
uses Java~21, Spring Boot 3.5.16, Spring Data JPA, H2, Flyway and the HAPI FHIR \emph{structures
and parser only} (\code{hapi-fhir-structures-r4} 8.12.1). HAPI server modules are excluded, because
they impose their own architecture and would remove the manipulated variable. Seed data is
synthetic (Synthea~\citep{walonoski2018synthea}).

We chose Spring Boot 3.5 rather than 4.x on purpose. Boot~4 (November 2025) changed core APIs
(Jackson~3, modularized starters) and is newer than much of the models' training data. It would
add a framework-novelty confound unrelated to architecture.

The variants differ only in code organization (ARCHITECTURES.md):
\begin{itemize}[itemsep=1pt]
  \item \NONE: idiomatic ``tutorial-style'' Spring. Controllers parse, validate, apply rules,
    map FHIR resources and call Spring Data repositories directly. Packages \code{controller},
    \code{entity}, \code{repository}, \code{util}. Realistic, not deliberately bad.
  \item \LAY: package-by-layer \code{web} $\rightarrow$ \code{service} $\rightarrow$
    \code{repository} (+ \code{entity}). HAPI types are confined to \code{web}.
  \item \HEX: ports and adapters after Hombergs~\citep{hombergs2023clean} (reference
    implementation \emph{buckpal}). A pure-Java \code{domain} with its own clinical model,
    \code{application} with inbound/outbound ports and use-case services, and
    \code{adapter.in.web}, \code{adapter.out.persistence} and \code{config}. FHIR mapping lives
    only in the web adapter.
\end{itemize}
The FHIR domain makes the placement of the standard model itself an architectural decision. In
\NONE{} HAPI resource classes run through the whole code base. In \LAY{} they stop at the service
boundary. In \HEX{} the domain has its own model and the FHIR mapping sits in one adapter, which is
exactly the kind of indirection that was costly in the predecessor study.

\paragraph{Equivalence gate.} A variant was admitted only if (i)~it passes the shared black-box
acceptance suite (79 tests over HTTP against the running jar, with FHIR R4 validation of every
response), (ii)~its Flyway schema is byte-identical and its dependency set exactly matches the
other variants (ArchUnit in test scope is the only extra), (iii)~it has zero violations of its own
ArchUnit rules, and (iv)~its size profile is published (Table~\ref{tab:size}). The size
differences are intrinsic to the styles (ports, mapping, a separate domain model). We report them
and do not equalize them. Each variant also contains an \code{AGENTS.md} with the same build/test
commands and a neutral, package-level layout description of equal length ($\pm$10\%) and no task
hints.

\begin{table}[t]\centering
\caption{Size profile of the three baselines at the equivalence gate (2026-10-07). All pass the
same 79-test acceptance suite.}
\label{tab:size}
\small
\begin{tabular}{lrrrrr}
\toprule
Variant & Main LOC & \code{.java} files & Packages & Own tests (incl.\ ArchUnit) & \code{AGENTS.md} words\\
\midrule
\NONE & 2243 & 27 & 4 & 8 & 173\\
\LAY & 2677 & 55 & 5 & 14 & 183\\
\HEX & 4064 & 104 & 12 & 16 & 185\\
\bottomrule
\end{tabular}
\end{table}

\paragraph{Construction.} The specification (\code{fhir/SPEC.md}) is the single source of truth
for behavior. Three AI agents built the three variants in parallel from it, and the author
reviewed them. Whenever two builders interpreted the specification differently, the ambiguity
became a binding clarification with its own acceptance test, and all variants were re-gated. This
produced 13 clarifications (C1--C13), for example the order of checks on \code{PUT} (parse error
before id mismatch before unknown id before business rules), the status for unsupported HTTP
methods on known types (405), and the handling of codes that are valid in FHIR R4 but outside the
supported subset. The gate caught four real behavioral divergences (C1, C7, C9, C11). We see this
as a side result: parallel independent implementation is an effective way to find specification
gaps.

\subsection{Tasks}\label{sec:tasks}

There are 24 tasks, written as product tickets. They describe externally observable behavior and
never name files, classes, packages or layers, so all three variants receive identical text
(1861 words in total). HTTP contracts and configuration property names are part of the external
interface and may be named. Every ticket ends with acceptance criteria that the agent can see. The
tasks fall into four strata of six, interleaved in a fixed order that mimics product evolution
(Table~\ref{tab:tasks}).

\begin{table}[t]\centering
\caption{Task strata. Units = architectural units a change is expected to touch.}
\label{tab:tasks}
\small
\begin{tabularx}{\textwidth}{@{}lp{2.2cm}Xl@{}}
\toprule
Stratum & Tasks & Example & Units\\
\midrule
D domain rule & T01 T05 T09 T13 T17 T21 & T01: body-weight Observations must use UCUM \code{kg} or \code{g}, else 422 \code{business-rule} & 1\\
F feature & T02 T06 T10 T14 T18 T22 & T02: \code{Observation?date=} search with prefixes; listed in the CapabilityStatement & 3+\\
C cross-cutting & T03 T07 T11 T15 T19 T23 & T03: optimistic locking with \code{If-Match} on every \code{PUT}; 412 on conflict & many\\
I integration & T04 T08 T12 T16 T20 T24 & T04: POST critical glucose Observations to a configured URL; failure must not fail the create & new adapter\\
\bottomrule
\end{tabularx}
\end{table}

\paragraph{Hidden tests.} Each task has hidden black-box acceptance tests (141 test methods in total,
3--8 per task), the same file for all variants. They run over HTTP against the running application.
FAIL\_TO\_PASS tests fail on every baseline. Guard tests (PASS\_TO\_PASS) pass on the baseline and stop
over-reaching solutions. For integration tasks the grader starts the application with the relevant
property pointing to a mock server controlled by the test. \code{harness/validate\_tasks.py}
checked on all three baselines that each task's FAIL\_TO\_PASS tests fail and the baseline suite
passes.
\TODO{attach the 24-task validation log; data/validate\_phase1.log covers 12 tasks only. No
reference solutions are in the repository; do not claim them.}

\paragraph{Dependencies and superseded tests.} Tasks were designed to be independent: 22 have no
dependency, and two have a soft one that their hidden tests do not need (T19 on T03, T20 on T06).
Some tickets legitimately change baseline behavior. For example, T18 makes \code{\_count} a valid
search parameter, and several tasks extend the CapabilityStatement that a baseline test checks
exactly. Nine of the 24 tasks therefore declare which baseline tests they \emph{supersede}. From
that task on, those tests are no longer counted as regressions (\S\ref{sec:harness}).

\subsection{Agent and models}\label{sec:agent}

\paragraph{Scaffold.} mini-swe-agent 2.4.6~\citep{minisweagent} with its \code{DefaultAgent}
and a single bash tool through native tool calling. The prompt is identical in every condition and
neutral about architecture. It asks the agent to read \code{AGENTS.md}, find the relevant code,
make the change, build and test, and submit. The agent may run the visible tests (\code{mvn test})
and start the application, and it never sees the hidden tests. The budget is 60 model calls per task.
\emph{Success} means solved within that budget: the hidden tests pass on the committed state,
whatever the stop reason. Each shell command times out after 600\,s. There is no wall-clock limit,
because provider latency must never become an agent failure.

\paragraph{Models.} Both arms use open-weights models, served through OpenRouter at temperature
0.2 with at most 16,384 output tokens per call:
\begin{itemize}[itemsep=1pt]
  \item \emph{strong}: Poolside Laguna~S~2.1~\citep{poolside2026laguna} (118B total / 8B active
    coding mixture-of-experts, OpenMDW-1.1 weights), run \emph{without} extended thinking. With
    thinking on, the serving provider ignored both the reasoning budget and \code{effort: low}. Calls
    ran to the output cap without a tool call, and in a shakedown 6 of 17 attempts ended in
    repeated format errors. With thinking off there were no format errors in 120 calls. This
    decision was made on reliability alone, before registration.
  \item \emph{weak}: Qwen3-Coder-30B-A3B-Instruct~\citep{qwen3report}, the model used in the
    cloud runs of the predecessor study. A preregistered floor rule would have replaced it with
    Qwen3-Coder-Next if it had solved fewer than 30\% of pilot attempts in \NONE. The rule was not
    triggered.
\end{itemize}
Open weights matter for reproducibility. Hosted endpoints change or disappear (\S\ref{sec:lessons}),
but anyone can re-run the study with the same weights. The serving provider is logged per call. In
the pilot, all strong-arm calls were served by Poolside and all weak-arm calls by SiliconFlow.

\subsection{Harness}\label{sec:harness}

\paragraph{Sandbox.} Each attempt runs in a fresh container of a pre-warmed image
(\code{archagents/sandbox:1}, Temurin JDK~21, Maven with every dependency cached) with
\code{--network none}, 2 CPUs and 3\,GB RAM. Maven runs offline. The trajectory's repository is
mounted as the workspace. Every command is wrapped in \code{timeout -k 5 590} inside the container,
so a hung build or a forgotten server is killed before the scaffold gives up on the command.

\paragraph{Grading.} After the agent stops, the harness commits the workspace. The grader exports the
committed tree (\code{git archive}), builds it in a separate offline container, starts the jar,
and runs the baseline acceptance suite plus the hidden tests of the current and all earlier tasks.
App configuration accumulates, so earlier integration tasks keep being checked against their mock
servers. A task is a \emph{success} if the snapshot builds, the application starts, and every
hidden test of the current task passes. A test is a \emph{regression} if it passed at the previous
grading and fails now. Tests of earlier tasks that the agent never solved are recorded as
\emph{still failing}, not as regressions. Superseded baseline tests are excluded from both.

\paragraph{Infrastructure failures.} Provider errors, container crashes and harness exceptions are
not agent failures. Only four exit states count as agent outcomes: submitted, step limit reached,
time exceeded, and repeated format errors. Any other exit rolls the workspace back to the last
commit, and the attempt is retried later, up to three times. After that it is recorded as missing.
A trajectory with more than 20\% infrastructure failures is re-run from scratch. No data are
excluded because of their outcome values.

\paragraph{Scheduling.} Task order is fixed. Execution is interleaved: for each task position, the
harness visits the trajectories in a seeded random order. Provider drift, throttling or a withdrawn
model therefore hits all architecture arms equally. As a logged deviation (\S\ref{sec:prereg}),
Phase~1 runs one such interleaving worker per run, two to three in parallel. Each worker contains
all six architecture $\times$ model trajectories of its run.

\subsection{Measures}\label{sec:measures}

\paragraph{Success and cost (per attempt).} Success (binary), regression flags, stop reason, total
tokens (prompt + completion, including cached prompt tokens), model calls, wall time, files
changed, and added/deleted lines. Total tokens is the primary cost measure. We also report a
recomputed USD cost that prices cached tokens separately.\footnote{Prices per million tokens at
registration (input / output / cache read): strong 0.09 / 0.18 / 0.009, weak 0.07 / 0.28 / 0.07 USD.}

\paragraph{Erosion (per committed snapshot).} \code{harness/erosion.py} measures every snapshot
post hoc in a second pinned image (\code{archagents/sandbox:2}). The image adds ArchUnit 1.5.1, CK
0.7.0, PMD/CPD 7.28.0 (\code{quickstart} rule set) and SpotBugs 4.10.4 (default effort, medium
threshold), all checked against published SHA-256 digests. Because a snapshot is deterministic,
post-hoc measurement gives the same values as measuring inside the loop. CK 0.7.0 cannot parse
Java records. The tool therefore parses at JLS~21 and runs CK on a copy in which each record is
rewritten as a final class with the same fields. With this change, CK reports no errors on any
baseline. The index has seven components (Table~\ref{tab:erosion}).

\begin{table}[t]\centering
\caption{Components of the erosion index (PREREGISTRATION \S4.3).}
\label{tab:erosion}
\small
\begin{tabular}{@{}rll@{}}
\toprule
$m$ & Component & Source\\
\midrule
1 & package-level dependency cycles in \code{org.archagents.fhir..} & ArchUnit slices\\
2, 3 & mean and max coupling between objects (CBO) & CK\\
4 & mean lack of cohesion (LCOM) & CK\\
5 & mean normalized distance from the main sequence $D=|A+I-1|$ & ArchUnit metrics\\
6 & static-analysis warnings (PMD violations + SpotBugs bugs) & PMD, SpotBugs\\
7 & duplicated lines (clones $\ge$100 tokens) & CPD\\
\bottomrule
\end{tabular}
\end{table}

For trajectory $t$, task index $k$ (0 = baseline) and component $m$, let $x_{t,k,m}$ be the metric
value on the committed snapshot. Then
\begin{align}
d_{t,k,m} &= x_{t,k,m} - x_{t,0,m}, &
s_m &= \operatorname{SD}_t\!\left(d_{t,K_t,m}\right), &
E_{t,k} &= \frac{1}{|M^\ast|}\sum_{m\in M^\ast}\frac{d_{t,k,m}}{s_m},
\end{align}
where $K_t$ is the last compiling snapshot of $t$, the standard deviation pools all 30 trajectories
across arms and models, and $M^\ast=\{m : s_m>0\}$. A component with zero pooled spread is dropped
and reported. There is no centering, so $E=0$ still means ``no change since baseline'' (centering
would add a constant per component and not change any slope). The H2 outcome is the OLS slope
$\beta_t$ of $E_{t,k}$ over $k$. Snapshots that do not compile have no metric values. They are left
out of the slope fits, and their number per arm is reported.

\paragraph{Conformance (H6).} For \LAY{} and \HEX{}, the number of violations of the variant's own
layer or hexagon rules (ARCHITECTURES.md). They are evaluated by a frozen copy of the baseline rules
in the measurement tool, never by the copy in the snapshot, which the agent may have edited. The
general checks (field injection, HAPI server modules) and the cycle rule are not counted, since
cycles enter the erosion index. A rule whose target package was deleted counts as one violation.
The H6 outcome is the OLS slope of this count over $k$.

\subsection{Hypotheses and analysis}\label{sec:analysis}

Table~\ref{tab:hyp} lists the hypotheses and their tests. All primary tests are permutation tests
on the unit of inference~\citep{good2005permutation}, implemented in \code{analysis/analyze.py},
which was frozen with the registration and dry-run on the pilot data. Everything is seeded
(20261012).

\begin{table}[t]\centering
\caption{Preregistered hypotheses for Phase 1 and their primary tests. Contrasts: \LAY$-$\NONE,
\HEX$-$\NONE, \HEX$-$\LAY, Holm-corrected within each hypothesis. H7 (domain generality) needs
Phase~2.}
\label{tab:hyp}
\small
\begin{tabularx}{\textwidth}{@{}l>{\raggedright\arraybackslash}p{4.6cm}l>{\raggedright\arraybackslash}X@{}}
\toprule
& Prediction & Sides & Primary test\\
\midrule
H1 & cost: \NONE{} $<$ \LAY{} $<$ \HEX{} (total tokens) & directional & sign-flip test on per-task differences of mean $\log$ tokens (24 tasks)\\
H2 & erosion slope lower in \LAY{} and \HEX{} than in \NONE{} & directional & label permutation of per-trajectory index slopes, stratified by model\\
H3 & success rate differs between architectures & two-sided & sign-flip test on per-task success-rate differences\\
H4 & \HEX$-$\NONE{} success gap larger for the weak model & directional & sign-flip test on per-task interaction contrasts\\
H5 & \HEX{} cost penalty largest for multi-unit strata (F, C, I), smallest for D & directional & per-stratum H1 contrasts; permutation of the (F,C,I) vs.\ D contrast\\
H6 & conformance violations accumulate (slope $>0$); \LAY{} vs.\ \HEX{} two-sided & mixed & as H2 on violation-count slopes\\
\bottomrule
\end{tabularx}
\end{table}

\paragraph{Task-level tests (H1, H3, H4, H5).} The task is the unit of inference. For each task and
arm we pool all attempts over runs and models and compute the success rate (H3) or the mean
$\log(\text{total tokens}+1)$ (H1). For each contrast, the 24 per-task differences enter a two-sided
sign-flip permutation test on their mean, with 100,000 random sign patterns including the observed
one. When there are 16 or fewer tasks, as in the per-stratum H5 contrasts, all patterns are
enumerated. H4 applies the same test to per-task interaction contrasts (weak-model effect minus
strong-model effect, in percentage points). H5 permutes stratum labels to test whether the
\HEX$-$\NONE{} cost difference is larger in strata F, C and I than in D. Secondary H1 analyses
repeat the test on successful attempts only, model calls, files changed, wall time and recomputed
USD.

\paragraph{Trajectory-level tests (H2, H6).} We fit one OLS slope per trajectory and compare arms with
a two-sided permutation test that shuffles architecture labels within model strata (10 trajectories
per arm, 10,000 relabelings).

\paragraph{Inference criteria.} $\alpha=0.05$. The three architecture contrasts within each
hypothesis are Holm-corrected~\citep{holm1979simple}. A directional hypothesis counts as supported
only if the corrected test is significant \emph{and} the effect has the predicted sign. All effects
come with 95\% percentile CIs from a cluster bootstrap over tasks with 10,000
resamples~\citep{efron1993bootstrap}. Secondary mixed models estimate adjusted effects:
a logistic GLMM \code{success \textasciitilde{} architecture * model + (1|task) + (1|trajectory)},
LMMs with the same structure for $\log$ tokens, and \code{erosion\_index \textasciitilde{}
architecture * task\_index + (task\_index|trajectory)}. A sensitivity analysis adds
\code{dependency\_failed} to the success model. Where models and primary tests disagree, both are
reported. \TODO{GLMM/LMM and the dependency\_failed sensitivity analysis are not implemented in
analyze.py; implement exactly as preregistered and label any choices not fixed there.}

\paragraph{Sample size and sensitivity.} The design is 3 architectures $\times$ 2 models $\times$ 5
runs = 30 trajectories of 24 tasks, i.e.\ 720 attempts. The size was set by resources: the study is
self-funded with an API budget of about USD~30. A simulation (\code{analysis/power\_sim.py};
logistic model with task and trajectory random effects calibrated to the pilot) gave the primary
H3 test with 24 tasks $\times$ 5 runs a power of 0.93 at $\alpha=0.05$ for a 15 percentage-point
difference, and 0.86 at the Holm-adjusted $\alpha=0.0167$. For 10 points the power is 0.67 and 0.53.
Smaller differences are reported as inconclusive, not absent. The simulation also showed that
adding tasks raises power more than adding runs, because the task is the unit of inference. This is
why we chose 24 tasks $\times$ 5 runs over 12 $\times$ 10 at equal cost.
\TODO{the committed power\_sim.py has N\_TASKS = 12; archive the 24-task run that produced these
numbers.} A preregistered, outcome-blind cost rule would drop run R5 if the forecast spend after
R1 exceeded USD~30. Outcome: \RESULT{cost-rule-outcome}.

\subsection{Preregistration and open science}\label{sec:prereg}

Preregistration fixes the confirmatory analysis before the data exist and so separates it from exploration~\citep{nosek2018preregistration}. The open-ended idea was registered on OSF on 2026-10-05 (\url{https://osf.io/w467e}). The full
protocol was registered on 2026-10-13 at 09:25 UTC with the OSF Preregistration template
(\url{https://osf.io/r5t3d}), under embargo until 2027-07-31. The registered materials are git tag
\code{prereg-v1}, commit \code{3d347cb}. Before registration only harness-debugging data existed:
three pilot tasks P1--P3 in all six cells (18 attempts) and one shakedown trajectory over T01--T24
(\NONE, strong model). Both are excluded from every analysis. The models, settings, sandbox
images, tasks, hidden tests and analysis script did not change after the freeze. One deviation is
logged (DEVIATIONS.md): execution runs in two to three parallel workers, one per run, instead of a
single queue over all 30 trajectories, because a sequential queue would take about 70 hours. Every
worker interleaves all arms, so the protection against drift is kept. Further deviations:
\RESULT{deviations-list}.

% =====================================================================================
\section{Results}\label{sec:results}

\TODO{Fill from data/phase1/analysis.json after data collection ends. Report every preregistered
test, including non-significant ones.}

\subsection{Data collection}\label{sec:res-collection}

Data were collected from \RESULT{collection-start} to \RESULT{collection-end}
(Table~\ref{tab:collection}).

\begin{table}[t]\centering
\caption{Data collection summary (top) and task success rate by architecture and model, pooled
over runs and tasks (bottom).}
\label{tab:collection}\label{tab:success}
\footnotesize
\begin{tabular}[t]{@{}lr@{}}
\toprule
Runs completed (cost rule) & \RESULT{runs-final}\\
Trajectories / attempts & \RESULT{N-trajectories} / \RESULT{N-attempts}\\
Missing after 3 infra retries & \RESULT{N-missing}\\
Infrastructure retries & \RESULT{N-infra-retries}\\
Trajectories re-run ($>$20\% infra) & \RESULT{N-traj-rerun}\\
Submitted / step limit / format err. & \RESULT{stop-reasons}\\
Non-compiling snapshots N/L/H & \RESULT{noncompiling-by-arm}\\
Providers (strong / weak) & \RESULT{providers}\\
Total cost; per attempt & \RESULT{cost-total}; \RESULT{cost-per-attempt}\\
\bottomrule
\end{tabular}\\[6pt]
\begin{tabular}[t]{@{}lccc@{}}
\toprule
 & \NONE & \LAY & \HEX\\
\midrule
strong & \RESULT{succ-none-strong} & \RESULT{succ-lay-strong} & \RESULT{succ-hex-strong}\\
weak & \RESULT{succ-none-weak} & \RESULT{succ-lay-weak} & \RESULT{succ-hex-weak}\\
pooled & \RESULT{succ-none} & \RESULT{succ-lay} & \RESULT{succ-hex}\\
\bottomrule
\end{tabular}
\end{table}

\begin{table}[t]\centering
\caption{Task-level primary tests (24 tasks): mean per-task difference, 95\% cluster bootstrap CI,
sign-flip $p$, Holm-adjusted $p$. H3 in percentage points of success; H1 in $\log$ total tokens
(positive = first arm more expensive); H4 is the interaction contrast weak-model effect minus
strong-model effect in pp (secondary, no Holm correction; prediction: negative for \HEX$-$\NONE).}
\label{tab:task-tests}\label{tab:h3}\label{tab:h1}\label{tab:h4}
\small
\begin{tabular}{@{}lllll@{}}
\toprule
Contrast & Estimate & 95\% CI & $p$ & $p_\text{Holm}$\\
\midrule
\multicolumn{5}{@{}l}{\emph{H3 success (pp), two-sided}}\\
\crow{H3}{lay-none}{\LAY$-$\NONE}{pp}
\crow{H3}{hex-none}{\HEX$-$\NONE}{pp}
\crow{H3}{hex-lay}{\HEX$-$\LAY}{pp}
\midrule
\multicolumn{5}{@{}l}{\emph{H1 cost ($\Delta\log$ tokens, all attempts), directional}}\\
\crow{H1}{lay-none}{\LAY$-$\NONE}{est}
\crow{H1}{hex-none}{\HEX$-$\NONE}{est}
\crow{H1}{hex-lay}{\HEX$-$\LAY}{est}
\midrule
\multicolumn{5}{@{}l}{\emph{H4 capability interaction (pp), directional}}\\
\LAY$-$\NONE & \RESULT{H4-lay-none-pp} & \RESULT{H4-lay-none-ci} & \RESULT{H4-lay-none-p} & n/a\\
\HEX$-$\NONE & \RESULT{H4-hex-none-pp} & \RESULT{H4-hex-none-ci} & \RESULT{H4-hex-none-p} & n/a\\
\HEX$-$\LAY & \RESULT{H4-hex-lay-pp} & \RESULT{H4-hex-lay-ci} & \RESULT{H4-hex-lay-p} & n/a\\
\bottomrule
\end{tabular}
\end{table}

\subsection{H3: Success}\label{sec:res-h3}

Table~\ref{tab:success} gives success rates per cell and Table~\ref{tab:task-tests} the primary
test. H3 is \RESULT{H3-verdict}. The GLMM odds ratio for \HEX{} vs.\ \NONE{} is
\RESULT{H3-glmm-hex-none} (profile CI). Adding \code{dependency\_failed}
\RESULT{H3-depfailed-sensitivity}. Figure~\ref{fig:success-cost}a shows success per task.

\begin{figure}[t]\centering
\begin{tabular}{@{}cc@{}}
\maybefig{figures/success_by_task.pdf}{0.46\textwidth} &
\maybefig{figures/cost_by_arm.pdf}{0.46\textwidth}\\
\small (a) success per task and architecture, by model & \small (b) total tokens per attempt (log scale)
\end{tabular}
\caption{Success and cost by architecture. Placeholders until unblinding.}
\label{fig:success-cost}
\end{figure}

\subsection{H1: Cost}\label{sec:res-h1}

H1 is \RESULT{H1-verdict} (Table~\ref{tab:task-tests}, Figure~\ref{fig:success-cost}b). The share
of attempts that hit the 60-call budget was \RESULT{budget-hit-rate-by-arm} (\NONE{} / \LAY{} /
\HEX), which bounds how far cost is censored (\S\ref{sec:threats}). The LMM estimate for \HEX{}
vs.\ \NONE{} is \RESULT{H1-lmm-hex-none}. Secondary measures are in Table~\ref{tab:h1sec}.

\begin{table}[t]\centering
\caption{H1 secondary cost measures and H5 (cost by stratum). Cells: mean per-task difference
[95\% CI], $p$. H5 strata have six tasks each (exact sign-flip test; smallest attainable two-sided
$p=0.031$).}
\label{tab:h1sec}\label{tab:h5}
\footnotesize
\begin{tabular}{@{}llll@{}}
\toprule
\multicolumn{4}{@{}l}{\emph{H1 secondary measures (Holm-adjusted $p$)}}\\
Measure & \LAY$-$\NONE & \HEX$-$\NONE & \HEX$-$\LAY\\
\midrule
$\log$ tokens, successful attempts & \RESULT{H1sec-succtok-lay-none} & \RESULT{H1sec-succtok-hex-none} & \RESULT{H1sec-succtok-hex-lay}\\
model calls & \RESULT{H1sec-calls-lay-none} & \RESULT{H1sec-calls-hex-none} & \RESULT{H1sec-calls-hex-lay}\\
files modified & \RESULT{H1sec-files-lay-none} & \RESULT{H1sec-files-hex-none} & \RESULT{H1sec-files-hex-lay}\\
wall time (s) & \RESULT{H1sec-wall-lay-none} & \RESULT{H1sec-wall-hex-none} & \RESULT{H1sec-wall-hex-lay}\\
recomputed USD & \RESULT{H1sec-usd-lay-none} & \RESULT{H1sec-usd-hex-none} & \RESULT{H1sec-usd-hex-lay}\\
\midrule
\multicolumn{4}{@{}l}{\emph{H5: \HEX$-$\NONE{} $\Delta\log$ tokens by stratum}}\\
Stratum & Estimate & 95\% CI & $p$\\
\midrule
D domain rule & \RESULT{H5-D-est} & \RESULT{H5-D-ci} & \RESULT{H5-D-p}\\
F feature & \RESULT{H5-F-est} & \RESULT{H5-F-ci} & \RESULT{H5-F-p}\\
C cross-cutting & \RESULT{H5-C-est} & \RESULT{H5-C-ci} & \RESULT{H5-C-p}\\
I integration & \RESULT{H5-I-est} & \RESULT{H5-I-ci} & \RESULT{H5-I-p}\\
(F, C, I) $-$ D & \RESULT{H5-multi-minus-D} & & \RESULT{H5-multi-minus-D-p}\\
\bottomrule
\end{tabular}
\end{table}

\subsection{H5: Cost by task stratum}\label{sec:res-h5}

H5 is \RESULT{H5-verdict} (Table~\ref{tab:h5}).

\subsection{H4: Capability interaction}\label{sec:res-h4}

H4 is \RESULT{H4-verdict} (Table~\ref{tab:task-tests}).

\subsection{H2: Erosion}\label{sec:res-h2}

The index used the components \RESULT{H2-metrics-used}. Components dropped for zero pooled spread:
\RESULT{H2-metrics-dropped}. Table~\ref{tab:traj-tests} compares slopes, and
Figure~\ref{fig:erosion}a shows the index over task position. H2 is \RESULT{H2-verdict}.
Per-component slopes are given in the supplementary material (\RESULT{H2-component-table}). LMM architecture $\times$ task-index
interaction: \RESULT{H2-lmm-interaction}.

\begin{table}[t]\centering
\caption{Trajectory-level tests: mean per-trajectory OLS slope per task for each arm of the
contrast, and label-permutation tests stratified by model (10 vs.\ 10 trajectories). H2: erosion
index. H6: own-architecture conformance violations.}
\label{tab:traj-tests}\label{tab:h2}
\small
\begin{tabular}{@{}lllll@{}}
\toprule
Contrast & slope (first arm) & slope (second arm) & $p$ & $p_\text{Holm}$\\
\midrule
\multicolumn{5}{@{}l}{\emph{H2 erosion index, directional}}\\
\LAY$-$\NONE & \RESULT{H2-lay-none-a} & \RESULT{H2-lay-none-b} & \RESULT{H2-lay-none-p} & \RESULT{H2-lay-none-pholm}\\
\HEX$-$\NONE & \RESULT{H2-hex-none-a} & \RESULT{H2-hex-none-b} & \RESULT{H2-hex-none-p} & \RESULT{H2-hex-none-pholm}\\
\HEX$-$\LAY & \RESULT{H2-hex-lay-a} & \RESULT{H2-hex-lay-b} & \RESULT{H2-hex-lay-p} & \RESULT{H2-hex-lay-pholm}\\
\midrule
\multicolumn{5}{@{}l}{\emph{H6 conformance violations, two-sided}}\\
\HEX$-$\LAY & \RESULT{H6-hex-slope} & \RESULT{H6-lay-slope} & \RESULT{H6-hex-lay-p} & n/a\\
\bottomrule
\end{tabular}
\end{table}

\begin{figure}[t]\centering
\begin{tabular}{@{}cc@{}}
\maybefig{figures/erosion_index.pdf}{0.46\textwidth} &
\maybefig{figures/conformance.pdf}{0.46\textwidth}\\
\small (a) erosion index $E_{t,k}$, mean and 95\% band & \small (b) own-architecture violations (\LAY, \HEX)
\end{tabular}
\caption{Erosion and conformance over task position, per architecture and model. Placeholders
until unblinding.}
\label{fig:erosion}\label{fig:conformance}
\end{figure}

\subsection{H6: Conformance}\label{sec:res-h6}

Own-architecture violations changed by \RESULT{H6-lay-slope} per task in \LAY{} and
\RESULT{H6-hex-slope} in \HEX{} (difference \RESULT{H6-hex-lay-diff}, $p=$
\RESULT{H6-hex-lay-p}; Figure~\ref{fig:erosion}b). H6 is \RESULT{H6-verdict}. The rules violated
most often were \RESULT{H6-top-rules}.
\TODO{analyze.py tests only the HEX-vs-LAY slope difference for H6; add the preregistered
one-sample test of slope $>0$ per arm, labeled as such.}

\subsection{Exploratory analyses}\label{sec:res-explore}

These analyses were preregistered as exploratory. \emph{Regressions:}
\RESULT{regression-rate-by-arm} of graded attempts broke at least one previously passing test.
\emph{Localization:} \RESULT{localization-summary} (files read before the first edit vs.\ files
where the change landed). \emph{Qualitative coding} of a stratified sample of failed attempts:
\RESULT{qualitative-summary}.

% =====================================================================================
\section{Discussion}\label{sec:discussion}

\TODO{Write after unblinding. Keep to what the preregistered tests support. The skeleton below
lists what each outcome pattern would need to address. No claims are made here.}

\paragraph{If architecture protects (lower erosion slope in \LAY/\HEX, success not lower).}
\begin{itemize}[itemsep=1pt]
  \item Which erosion components carry the effect: cycles and coupling, or warnings and
    duplication? The baselines differ in size and structure, so a difference in components that
    scale with code size needs care.
  \item Is the protection bought with cost (H1)? Report the price per task and per point of
    erosion avoided.
  \item Did conformance hold (H6 slope near zero), or did agents keep the metrics flat while
    breaking layer rules?
  \item Relation to \citep{archequalizer2026,hexdebt2025}: does the code's structure act as an
    implicit specification?
\end{itemize}

\paragraph{If the agent breaks the architecture (H6 slope $>0$, H2 not supported).}
\begin{itemize}[itemsep=1pt]
  \item Which rules break first (e.g.\ HAPI types leaking into \code{domain}/\code{service},
    adapters calling services directly)? Link to the qualitative coding.
  \item Is the structural cost of \HEX{} still paid (H1) while its benefits erode? That is the
    worst case for practice.
  \item Interaction with capability (H4) and stratum (H5): do cross-cutting and integration tasks
    drive the violations?
  \item Implications for guardrails (enforced ArchUnit rules, instructions). This motivates
    Study~2.
\end{itemize}

\paragraph{If results are mixed or inconclusive.}
\begin{itemize}[itemsep=1pt]
  \item Use the sensitivity analysis: effects below about 10--15 pp are outside what this
    design can detect.
  \item Compare with the predecessor study's single-task tax, and say whether cumulative,
    larger-scale work changes the picture.
\end{itemize}

% =====================================================================================
\section{Threats to Validity}\label{sec:threats}

We follow the usual classification~\citep{wohlin2012experimentation}. The first items were stated
in the preregistration (\S7.1). The rest were identified while running the study.

\paragraph{Construct validity.}
\emph{Size is part of the treatment.} \HEX{} has 1.81$\times$ the LOC and 3.9$\times$ the files of
\NONE. We report the profile and do not equalize it, so effects of architecture cannot be
separated from the code volume the style implies. \emph{Erosion index.} The index aggregates
established metrics with equal weight. Some components, such as distance from the main sequence,
depend on package structure and may behave differently by style. We therefore report every
component separately. \emph{Conformance} exists only for \LAY{} and \HEX. \emph{Cost} includes
cached prompt tokens, which are cheap in money but still measure how much context the agent
processes. We report recomputed USD as a secondary measure. \emph{Step-limit censoring.} Cost is
right-censored at 60 calls. In the shakedown the strong model used the full budget in 18 of 24
tasks. Censoring lowers the power of H1 and compresses differences between arms. It does not bias
the permutation test under the null. The secondary analysis on successful attempts is only
partly a remedy, because success is itself an outcome. \emph{Excluded snapshots.} Non-compiling
snapshots are left out of the slope fits. If arms differ in compile rates, slopes describe
different subsets. We report the counts per arm.

\paragraph{Internal validity.}
\emph{Construction bias.} One author, assisted by AI agents, built all variants and tasks, and no
independent human review took place. Mitigations: a public reference style, the mechanical
equivalence gate (79 black-box tests, identical schema and dependencies), frozen ArchUnit rules,
published size profiles and full release of all materials. AI-built variants may also be closer
to what models find natural than human-written code would be. \emph{Task authoring.} Tickets are
behavior-only and graded by identical black-box tests, but nobody else reviewed them.
\emph{Provider drift.} Hosted models can change silently, and providers can re-route requests. We
pinned model IDs, interleaved all arms in every worker, and logged the provider per call. Drift
during the collection window affects all arms equally but could still interact with
architecture. \emph{Dependencies between tasks.} A failed task can make later tasks harder in that
trajectory. We report the preregistered \code{dependency\_failed} sensitivity analysis.

\paragraph{External validity.}
\emph{Single domain in Phase 1.} Results come from one application. FHIR is a public standard that
models have probably seen often, and its external model makes the placement of mapping code
especially salient. Phase~2 (order management) tests generality (H7). We claim nothing about other
languages, frameworks or much larger systems. \emph{Model choice.} Two open-weights models in one
scaffold, chosen partly for cost, with the strong model run without extended thinking.
Frontier proprietary agents with richer tools may behave differently, and with two models H4 is a
contrast between two points, not a capability curve. \emph{Scaffold.} A bash-only loop with a
60-call budget is deliberately minimal. Products with planning, retrieval or self-review steps
may close or widen any gap.

\paragraph{Conclusion validity.} 24 tasks give adequate power only for differences of about
15 pp or more (\S\ref{sec:analysis}). H5 per-stratum tests have six tasks each and are
descriptive. Contamination is limited by keeping all materials private until data collection
ends and by running without network access. The materials will become public afterwards.

% =====================================================================================
\section{Lessons for Running Agent Experiments on a Shoestring}\label{sec:lessons}

The study ran on one laptop with Docker and a budget of USD~30. Most of the engineering effort
went into problems that a well-funded lab may never see but that can silently bias a low-budget
experiment. We record them because each one would have turned infrastructure behavior into
apparent agent behavior.

\paragraph{1. Free endpoints are not infrastructure.} We first planned to use free
(\code{:free}) endpoints. \code{qwen/qwen3-coder:free} was withdrawn on 2026-10-11, and
\code{poolside/laguna-s-2.1:free} was throttled for many minutes per call. Both were dropped
before any confirmatory data in favor of the paid endpoints of the same models. With
provider-side prompt caching the paid runs were cheap: 83\% of prompt tokens in the pilot and
96--97\% in the shakedown were served from cache, and Phase~1 cost about USD~0.04 per attempt.
Open weights mean the study stays reproducible when endpoints change.

\paragraph{2. Providers may ignore reasoning settings.} With extended thinking enabled, the
provider ignored both \code{reasoning.max\_tokens} and \code{effort: low}. Calls ran to the output
cap without a tool call, and the scaffold counted them as format errors. Six of 17 shakedown
attempts ended this way, a failure mode unrelated to the code base that would have added noise to
every arm. Turning thinking off removed it (0 format errors in 120 calls). Measure the format-error
rate in a shakedown, and fix reasoning settings before registration on reliability grounds alone.

\paragraph{3. Errors can arrive with HTTP 200.} Some upstream errors came back as HTTP 200 with
an \code{error} body and no \code{choices}. The stock scaffold crashed on them with a
\code{KeyError}. A thin wrapper turns them into retryable API errors and raises the request
timeout from 60 to 300\,s. More generally, the harness keeps an explicit allow-list of exit states
that count as agent outcomes. Everything else is infrastructure: it is rolled back to the last
commit and retried, never scored.

\paragraph{4. Command timeouts do not kill processes.} When a command times out, the scaffold
stops waiting, but a hung \code{mvn test} or a server the agent started in the background can keep
running inside the container. It then competes for the 2-CPU budget and inflates the wall time of
later commands. Wrapping every command in \code{timeout -k 5 590} inside the container, 10\,s
before the scaffold's own 600\,s limit, and removing the container after each attempt fixed this.

\paragraph{5. Define ``regression'' with care.} In a cumulative design, the naive definition (any
earlier hidden test that fails now) counts every earlier task the agent never solved as a
regression in every later attempt. We define a regression as a test that passed at the previous
grading and fails now, record ``still failing'' tests separately, and compare against the last
snapshot that built and started.

\paragraph{6. Cumulative grading needs \emph{supersedes}.} Legitimate tickets change baseline
behavior: a new search parameter that used to be rejected, a stricter rule that an old test
violates, an extended CapabilityStatement that an old test checks exactly. Without an explicit
list of superseded tests, correct solutions would register as regressions. Nine of 24 tasks needed
such entries. A full shakedown trajectory over all 24 tasks, run before registration, found these
conflicts. Configuration must also accumulate, so that earlier integration tasks keep being tested
against their mock servers.

\paragraph{7. Make every step resumable and idempotent.} Long runs on a laptop get interrupted. If
the process is killed between writing a result row and saving trajectory state, a naive restart
grades the same task twice. The runner reconciles state with the results file at startup.
Measuring erosion post hoc from git snapshots, rather than in the loop, means a crashed or
upgraded measurement tool never costs agent runs.

\paragraph{8. Pin and adapt the measurement tools.} CK 0.7.0 predates Java records. It either mis-parsed
Java~21 syntax or lost every type in a file that contains a record (22 files in the \HEX{}
baseline, 3--4 in the others). Without the desugaring step (\S\ref{sec:measures}), coupling and
cohesion would have been computed on an incomplete set of types that differs by variant. We pinned tool versions in a separate image tag, checked release
digests, and verified that CK reports no errors on the baselines.

% =====================================================================================
\section{Conclusion}\label{sec:conclusion}

We set out to test, under preregistered control, whether layered and hexagonal architecture help
or hinder AI coding agents that evolve a Java service over many tasks, and whether the agents
preserve the architecture. \RESULT{conclusion-summary}. \TODO{Two to three sentences after
unblinding, limited to supported hypotheses; state Phase 2 and Study 2 (guardrails) as next steps.}

\paragraph{Data availability.} All materials are released under MIT (code) and CC-BY 4.0 (text,
data): the three variants, specification and clarifications, tasks, hidden tests, harness, agent
and model configuration, sandbox Dockerfiles, raw agent transcripts, per-task git snapshots,
results and erosion tables, and analysis scripts. All data are synthetic. OSF project:
\url{https://osf.io/rneym}. Registrations: \url{https://osf.io/r5t3d} (Phase~1, embargoed until
2027-07-31) and \url{https://osf.io/w467e} (idea). Code: \TODO{GitHub URL and Zenodo DOI}.

\paragraph{Use of AI tools.} AI coding agents built the three variants from the author's
specification under the author's review (\S\ref{sec:variants}). \TODO{Disclose AI assistance in
harness development and in drafting this manuscript, following the target venue's policy.}

\bibliographystyle{plainnat}
\bibliography{refs}

\end{document}
